\section{考察}

\subsection{手動計測と自動計測の比較と効率性()}
本実験では、実験1で手動計測を、実験3および4でLabVIEWを用いた自動計測を行った。実験1では13点のデータ取得に多くの時間を要したが、自動計測では高密度なサンプリングをそれよりも短い時間で完結させることができた。
特にダイオードの立ち上がり領域（$0.5\,\mathrm{V} \sim 0.7\,\mathrm{V}$付近）のような非線形性が強い区間において,
手動では困難な「細かな電圧ステップでの追従」が自動計測により容易となり,
図\ref{fig:graph_Ed_I}のような滑らかな特性曲線を得ることができた。
これは、急峻な変化を伴う電子部品の評価において,
コンピュータによる高速制御が不可欠であることを示唆している。

\subsection{LabVIEWの処理について}
LabVIEWが実行している処理の整理本実験の自動計測プログラム（\texttt{Lab4\_USB Auto Mes.vi}など）において、
LabVIEWがバックグラウンドで行っている処理は以下の通りである。
\begin{enumerate}
\item 各リソースとの通信確立と初期設定: VISAリソースを介してファンクションジェネレータ（FG）およびディジタルマルチメータ（DMM）との接続を確立し、測定レンジや出力モードを初期設定する。
\item ループの開始 : ブロックダイヤグラムの処理が始まり, FGブロックに数値が渡され始める。
\item コマンドの送信: FGに対して電圧設定コマンドを送信する。(LabVIEWからの送信をもとにファンクションジェネレータが電圧を回路に入力する)
\item データの取得と変換: DMMに対して測定要求コマンドを送信し、返ってきた計測文字列（ASCII）を数値データに変換する。
\item 演算処理: 2,3を繰り返して取得した抵抗電圧 $E_R$ から、ブロックダイヤグラムでプログラムされた回路に基づいて、電流 $I$ およびダイオード間電圧 $E_D$ を即座に算出する。
\item 可視化と保存: 算出した数値をフロントパネルのグラフにプロットし、終了後に一括してテキストファイル（.txt）へ書き出す。
\end{enumerate}

\subsection{ゲイン特性の自動計測について}
具体的な手法を以下に提案する。ただしファンクションジェネレータをFC, ディジタルマルチメータをDMM1, DMM2を表記する。
\begin{enumerate}
    \item  LRC直列回路を組み立てる。\\
    電圧入力をFCで,
    回路の入力電圧 $V_{in}$ をDMM1で,
    コンデンサの出力電圧 $V_{out}$ をDMM2で入出力できるように接続する。
    \item  LabVIEW上で以下の回路を組み立てる。\\
    \begin{itemize}
		\item 周波数を変化させるループ
    	\item FCブロックに周波数を入力, DMM1およびDMM2から電圧の実効値を出力するような回路
    	\item 電圧の実効値から$G = 20\times\log_{10}{V_{out} / V_{in}}$ を算出する回路
    	\item 横軸を周波数比$\omega / \omega_0$（対数軸, $\omega_0$は共振周波数）,縦軸をゲイン$G$としたボード線図を自動描画回路
	\end{itemize}
    \item  実行し, グラフを得る。
\end{enumerate}
これにより,
手動では膨大な時間を要する共振曲線のプロットを極めて短時間で,
かつ正確に行うことができる。
  
\subsection{ダイオードの材料・構造による特性の差異}
実験4の比較により、各ダイオードの材料・構造に起因する物理的知見および実用上の利点が以下のように整理された。

\begin{itemize}
    \item \textbf{シリコンダイオード：}
    順方向バイアスにおいて約$0.6\,\mathrm{V}$付近から電流が急峻に立ち上がる特性が確認された。ゲルマニウムに比べ逆方向漏れ電流が極めて小さく、熱安定性に優れるという特徴を持つ。
    \textbf{応用例：} 一般的な整流回路、電源保護回路、AC-DC変換アダプタなど。

    \item \textbf{ゲルマニウムダイオード：}
    シリコンよりも大幅に低い約$0.2\,\mathrm{V}$付近から電流が流れ始めることが確認された。これはゲルマニウムの禁制帯幅（バンドギャップ）が小さいためであり、微小信号への応答性に優れる。
    \textbf{応用例：} AMラジオの検波回路、低電圧動作が要求される高周波スイッチング回路など。

    \item \textbf{ツェナーダイオード：}
    順方向特性はシリコンと同様であったが、逆方向において約$-5.5\,\mathrm{V}$付近で急激に電流が流れる「ツェナー降伏」が観測された。降伏領域では電圧変化に対して電流が鋭敏に変化するため、端子電圧を一定に保つ性質がある。
    \textbf{応用例：} 定電圧源（基準電圧）の作成、過電圧保護回路、リミッタ回路など。
\end{itemize}
\subsection{手動測定の直列回路データからダイオードの特性を得る方法}
実験1で計測したのは抵抗とダイオードの直列回路による特性であるが, 
ここからダイオードのみの特性を求める場合, 実験3で計算したようにするとよい。
具体的には,表\ref{tab:diode_manual}で直流印加電圧$E$, 抵抗両端電圧$E_R$に対して, ダイオードに加わる電圧$E_D$は次式で求まる。
\begin{equation}
    E_D = E - E_R
\end{equation}
これはアノードカソード間電圧に相当する。
ダイオードを流れる電流は図\ref{tab:diode_manual}の電流$I$に同じ。
よって,$E_D$と$I$の関係を調べることでダイオード単体の静特性を調べることができる。
実際に調べたものが実験3の図\ref{fig:graph_Ed_I}である。
図\ref{fig:graph_E_I}では抵抗の電圧降下分が含まれるため立ち上がりが緩やかであるが, 
\ref{fig:graph_Ed_I}ではダイオードの非線形な特性が確認できる。


\section{結論}
本実験では、LabVIEWを用いた自動計測システムを構築し、各種ダイオードの電圧-電流静特性を評価した。
手動計測と比較して、自動計測は圧倒的に高い作業効率とデータ密度を実現できることを確認した。
また、VISA制御によるシステムが理論値通りの精度で動作することを検証し、
取得したデータからシリコン、ゲルマニウム、およびツェナーダイオードの物理的特性の差異を明確に示すことができた。
以上の結果から、コンピュータによる自動計測は、複雑な電子回路や素子の特性解析において極めて有効な手法であることを結論付ける。